\documentclass[11pt]{article}
\usepackage[margin=1in]{geometry}
\usepackage{booktabs}
\usepackage{amsmath}
\usepackage{siunitx}
\usepackage{caption}
\usepackage{hyperref}

\title{Weighting Scheme: Hierarchical Risk Parity (HRP)}
\author{MetaFVG Mean-Reversion Portfolio -- 14-Instrument Diversified Universe}
\date{\today}

\begin{document}
\maketitle

\section{Scheme Description}

The portfolio blends 14 instruments trading a Spearman-gated MetaFVG mean-reversion
strategy, deliberately assembled across different central banks and currency blocs
(AUDUSD, EURUSD, GBPUSD, NZDUSD, USDCAD, USDCHF, USDJPY, US500, GER40, JP225, HK50,
USDSGD, EURNOK, USDZAR). The current production blend uses equal weights
($w_i = 1/14$ for all $i$). This report tests \emph{Hierarchical Risk Parity} (HRP,
L\'{o}pez de Prado 2016, ``Building Diversified Portfolios that Outperform
Out-of-Sample''), which replaces the quadratic-optimization step of mean-variance and
risk-parity methods with a tree-structured allocation derived from the correlation
matrix's hierarchical clustering. HRP was designed to be robust to estimation error
in the covariance matrix (no matrix inversion) by allocating weight through a
dendrogram rather than solving a global optimization.

\subsection{Algorithm}

\textbf{1. Distance matrix.} Convert the correlation matrix $\rho$ into a proper
metric distance,
\[
  d_{ij} = \sqrt{\tfrac{1}{2}\left(1 - \rho_{ij}\right)}.
\]

\textbf{2. Hierarchical clustering.} Condense $d$ into its upper-triangular form and
apply single-linkage agglomerative clustering to build a dendrogram over the 14
instruments.

\textbf{3. Quasi-diagonalization.} Read the leaf order off the dendrogram and reorder
the instruments by it, so that instruments merged early in the tree (more correlated)
sit adjacent to one another in the resulting sequence.

\textbf{4. Recursive bisection.} Walk down the tree top-down. At each split of a
cluster $c$ into $c_0$ (first half of the quasi-diagonal order) and $c_1$ (second
half), each sub-cluster's variance is computed using \emph{inverse-variance
allocation within the sub-cluster} (diagonal of the covariance matrix only, HRP's
standard simplification -- no sub-cluster covariance terms):
\[
  \mathrm{IVP}_i = \frac{1/\Sigma_{ii}}{\sum_{k \in c} 1/\Sigma_{kk}}, \qquad
  V(c) = \mathrm{IVP}^\top \Sigma_c\, \mathrm{IVP}.
\]
The split factor allocates \emph{more} weight to the lower-variance side:
\[
  \alpha = 1 - \frac{V(c_0)}{V(c_0) + V(c_1)}, \qquad
  w_{c_0} \mathrel{*}= \alpha, \qquad w_{c_1} \mathrel{*}= (1-\alpha).
\]
This repeats recursively until every cluster is a singleton, giving the final weight
vector $w$ (renormalized to sum to 1).

\textit{Implementation note: all covariance/variance products above were computed via
elementwise (Hadamard) broadcast-and-sum, e.g.\ \texttt{(np.outer(ivp, ivp) *
sub).sum()}, rather than matrix multiplication, and the covariance matrix itself was
built from a hand-rolled pairwise-complete estimator rather than
\texttt{DataFrame.cov()} -- both \texttt{DataFrame.cov()} and \texttt{Series.cov()}
were found to crash this environment's BLAS/LAPACK backend (exit 127, no traceback),
in addition to the previously known \texttt{np.corrcoef} / matrix-matmul crash mode.
\texttt{DataFrame.corr()} itself is unaffected and was used directly for $\rho$.}

\section{Quasi-Diagonalized Cluster Order}

The dendrogram leaf order groups instruments by realized correlation structure:

\begin{center}
\texttt{EURNOK -- JP225 -- US500 -- GER40 -- GBPUSD -- EURUSD -- AUDUSD -- NZDUSD -- USDJPY -- HK50 -- USDCAD -- USDZAR -- USDCHF -- USDSGD}
\end{center}

Several groupings match fundamental intuition: the three equity indices \textbf{JP225,
US500, GER40} land adjacent despite trading different central-bank regimes and
sessions -- consistent with a shared global risk-on/risk-off factor dominating their
correlation. \textbf{AUDUSD} and \textbf{NZDUSD}, the classic commodity-bloc pair, are
adjacent as expected. \textbf{GBPUSD} and \textbf{EURUSD}, the two large European
majors, are adjacent as well. \textbf{EURNOK} anchors one end of the order (Norway's
oil-linked currency, least like the rest of the book) and \textbf{USDSGD} anchors the
other. \textbf{HK50} sits with \textbf{USDJPY} rather than with the other equity
indices -- a less intuitive pairing, likely an artifact of single-linkage clustering
``chaining'' through a weak intermediate correlation rather than a strong direct
relationship.

\section{HRP Weights}

\begin{table}[h]
\centering
\caption{HRP weights, sorted descending. Ratio column is $w_i / (1/14)$.}
\begin{tabular}{lS[table-format=1.5]S[table-format=1.3]}
\toprule
Symbol & {Weight} & {Weight / (1/14)} \\
\midrule
USDSGD & 0.20604 & 2.885 \\
USDCAD & 0.14479 & 2.027 \\
EURUSD & 0.12427 & 1.740 \\
EURNOK & 0.09692 & 1.357 \\
GBPUSD & 0.08697 & 1.218 \\
USDCHF & 0.07841 & 1.098 \\
USDJPY & 0.07236 & 1.013 \\
AUDUSD & 0.04818 & 0.675 \\
NZDUSD & 0.04477 & 0.627 \\
USDZAR & 0.02834 & 0.397 \\
US500  & 0.02246 & 0.314 \\
GER40  & 0.02223 & 0.311 \\
HK50   & 0.01393 & 0.195 \\
JP225  & 0.01033 & 0.145 \\
\midrule
Sum & 1.00000 & 14.000 \\
\bottomrule
\end{tabular}
\end{table}

HRP's recursive bisection pushes weight heavily toward the lower-variance FX legs --
USDSGD, USDCAD, EURUSD, and EURNOK all land at 1.4--2.9$\times$ equal weight -- and
away from the four equity indices, which are squeezed down to roughly 0.15--0.31$\times$
equal weight (JP225 and HK50 the most, at $\sim$0.15--0.20$\times$). This mirrors the
direction seen with ERC (equity indices underweighted), but is considerably more
extreme here because HRP's top-level bisection groups three of the four equity indices
together early in the quasi-diagonal order, and inverse-variance allocation then
compounds against that whole cluster relative to the large low-variance FX block on
the other side of the split.

\section{Correlation Impact}

\begin{table}[h]
\centering
\caption{Realized correlation and diversification impact, equal-weight vs.\ HRP.}
\begin{tabular}{lS[table-format=1.6]S[table-format=1.6]}
\toprule
Metric & {Equal Weight} & {HRP} \\
\midrule
Weighted avg.\ pairwise correlation & 0.015615 & 0.017022 \\
Diversification ratio (DR) & 2.080922 & 1.538298 \\
\bottomrule
\end{tabular}
\end{table}

HRP does not reduce realized correlation here -- the weighted average pairwise
correlation \emph{rises} from 0.0156 to 0.0170 (a $\sim$9\% relative increase), because
the scheme concentrates weight in a cluster of FX pairs (USD-quoted majors/crosses)
that are, on net, slightly \emph{more} mutually correlated than the full 14-instrument
book, while pulling weight away from the equity indices that anchor much of the
book's diversification. The diversification ratio also falls sharply (2.081 $\to$
1.538): DR rewards spreading weight across instruments with high standalone volatility
relative to blended portfolio volatility, and HRP's variance-minimizing objective
does the opposite by construction.

\section{Results}

\begin{table}[h]
\centering
\caption{Quantstats metrics, equal-weight vs.\ HRP blend.}
\begin{tabular}{lS[table-format=1.6]S[table-format=1.6]}
\toprule
Metric & {Equal Weight} & {HRP} \\
\midrule
Sharpe & 0.517941 & 0.169721 \\
Sortino & 0.816977 & 0.264785 \\
Calmar & 0.343017 & 0.106463 \\
CAGR & 0.000183 & 0.000057 \\
Max Drawdown (\%) & -5.325228 & -5.325228 \\
Volatility (\%) & 3.527622 & 3.343598 \\
Tail Ratio & 1.238689 & 1.002700 \\
Ulcer Index & 0.000192 & 0.000295 \\
Kelly Criterion & 0.066763 & 0.027229 \\
VaR (\%) & -0.358268 & -0.344199 \\
CVaR (\%) & -0.715011 & -0.797777 \\
Skew & 1.544839 & 1.976561 \\
Recovery Factor & 2.260123 & 0.701973 \\
\midrule
Win rate & 0.308850 & 0.298013 \\
Best day & 0.000291 & 0.000291 \\
Worst day & -0.000239 & -0.000239 \\
Total return & 0.001204 & 0.000374 \\
\bottomrule
\end{tabular}
\end{table}

\section{Takeaway}

HRP fails on every axis this exercise cares about. The weighted average pairwise
correlation \emph{increases} (0.0156 $\to$ 0.0170) instead of falling, the
diversification ratio drops by nearly a quarter (2.081 $\to$ 1.538), and Sharpe
collapses by two-thirds (0.518 $\to$ 0.170) alongside similarly large drops in
Sortino, Calmar, and recovery factor. The mechanism is clear from the weights table:
HRP's variance-based recursive bisection has no notion of expected return and
systematically defunds the equity-index legs -- JP225 down to 0.15$\times$ and HK50 to
0.20$\times$ equal weight -- purely because their standalone P\&L variance is higher,
even though those same instruments are exactly the ones providing this book's
fundamental (cross-asset-class) diversification. On this data, equal weighting
remains clearly superior; HRP should not replace it for this portfolio.

\end{document}
